\lysh{38845}{4}

\begin{alist}
\item Τα συνολικά έσοδα της εταιρείας ανά ημέρα είναι $25 \cdot 100 = 2500\text{\euro}$ και το συνολικό κόστος παραγωγής $5 \cdot 100 = 500\text{\euro}$. Άρα το κέρδος ανά ημέρα είναι $2500 - 500 = 2000\text{\euro}$.

\item Αν η εταιρεία μειώσει την τιμή της φανέλας κατά $2\text{\euro}$, οι συνολικές πωλήσεις αναμένεται να διαμορφωθούν στις $100 + 2 \cdot 10 = 120$ μονάδες την ημέρα. Οπότε, τα συνολικά έσοδα θα είναι $23 \cdot 120 = 2760\text{\euro}$ και το συνολικό κόστος $5 \cdot 120 = 600\text{\euro}$. Άρα το κέρδος θα είναι $2760 - 600 = 2160\text{\euro}$.

\item Αν η εταιρεία μειώσει την τιμή κατά $x\text{\euro}$, δηλαδή στα $25 - x\text{\euro}$ τη φανέλα, θα πουλήσει $100 + 10x$ φανέλες. Οπότε τα συνολικά έσοδα θα είναι $(25 - x)(100 + 10x)\text{\euro}$. Αντίστοιχα, το συνολικό κόστος θα είναι $5 \cdot (100 + 10x)\text{\euro}$. Άρα, το κέρδος $P$ δίνεται από τη σχέση
\begin{gather*}
P(x) = (25 - x)(100 + 10x) - 5 \cdot (100 + 10x) = \\
(100 + 10x)(20 - x) = -10x^2 + 100x + 2000\text{\euro}.
\end{gather*}

\item
\begin{rlist}
\item Από την καμπύλη της συνάρτησης του κέρδους προκύπτει ότι η εταιρεία μπορεί να μειώσει την τιμή κάθε φανέλας το πολύ κατά $20\text{\euro}$, καθώς για μεγαλύτερη μείωση το κέρδος $P$ παίρνει αρνητικές τιμές, που σημαίνει ότι η εταιρεία έχει ζημιά.

\textit{Άλλος τρόπος:}
Για να μην έχει ζημιά η εταιρεία, θα πρέπει η τιμή πώλησης κάθε φανέλας να μην είναι μικρότερη από το κόστος παραγωγής της. Άρα, η ελάχιστη τιμή πώλησης είναι $5\text{\euro}$, οπότε η εταιρεία μπορεί να μειώσει την τιμή της φανέλας το πολύ κατά $20\text{\euro}$.

\item Το κέρδος $P$ που έχει τώρα η εταιρεία αντιστοιχεί στην τιμή $x = 0$, και είναι $2000\text{\euro}$.
Από την καμπύλη της συνάρτησης κέρδους προκύπτει ότι για να μην έχει μικρότερο κέρδος από το τωρινό, η εταιρεία μπορεί να μειώσει την τιμή το πολύ κατά $10\text{\euro}$.
Άρα η μικρότερη τιμή με την οποία μπορεί να πουλάει τη φανέλα είναι $25 - 10 = 15\text{\euro}$.
\end{rlist}
\end{alist}
